\section{A concrete target beyond Dirichlet GRH}\label{sec:cubic}

\subsection{The field and the missing factor}

Let
\[
F=\Q(\theta),\qquad \theta^3-\theta-1=0,
\qquad K=\Q(\sqrt{-23}),
\]
and let \(H\) be the splitting field of the cubic. The following identities are classical; the example and its factorization are described in Venkatesh's lectures, in notes by Feng \cite{venkatesh}, and the Hilbert-class-field identification also occurs in Schoof's paper \cite{schoof}.

\begin{proposition}\label{prop:cubic}
The group \(\Cl(K)\) is cyclic of order \(3\), and \(H\) is its Hilbert class field. If \(\psi\) is either nontrivial character of \(\Cl(K)\), then
\begin{align}
\zeta_F(s)&=\zeta(s)L_K(s,\psi),\label{eq:cubic}\\
\zeta_H(s)&=\zeta_K(s)L_K(s,\psi)^2
=\zeta(s)L(s,\chi_{-23})L_K(s,\psi)^2.\label{eq:galois}
\end{align}
\end{proposition}

\begin{proof}
The only possible rational roots of \(X^3-X-1\) are \(\pm1\), and neither is a root. Its discriminant is
\[
-4(-1)^3-27(-1)^2=-23.
\]
It is squarefree, so the order \(\Z[\theta]\) is maximal; it is nonsquare, so the transitive Galois group is \(S_3\), with quadratic subfield \(K\).

Reduced positive definite primitive forms of discriminant \(-23\) have first coefficient at most \(\sqrt{23/3}<3\). Enumerating the possibilities gives exactly
\[
(1,1,6),\qquad (2,1,3),\qquad (2,-1,3).
\]
The classical ideal-class/form correspondence therefore gives \(h_K=3\).

The extension \(H/K\) is cyclic of degree \(3\). No prime other than \(23\) ramifies. Modulo \(23\),
\[
X^3-X-1=(X-10)^2(X-3).
\]
The simple root lifts over \(\Q_{23}\), so the local decomposition group fixes a root and is contained in an order-two stabilizer. The quadratic field is ramified at \(23\), so inertia is exactly that order-two subgroup; its intersection with \(A_3=\Gal(H/K)\) is trivial. Thus \(H/K\) is unramified at finite places. There are no real places of \(K\) to consider. Its degree is the class number, so class field theory identifies it with the Hilbert class field.

Let \(V\) be the standard two-dimensional representation of \(S_3\). The permutation representation on three roots is \(\mathbf1\oplus V\); hence its Artin \(L\)-function gives \(\zeta_F=\zeta L_\Q(V)\). Inducing either nontrivial character of \(A_3\) gives \(V\): the resulting character takes values \(2,0,-1\) on the identity, transpositions, and three-cycles. Compatibility of Artin \(L\)-functions with induction identifies \(L_\Q(V)\) with \(L_K(\psi)\). The two nontrivial characters induce isomorphic representations and give the same \(L\)-function. Finally the regular representation of \(S_3\) is \(\mathbf1\oplus\mathrm{sgn}\oplus V^{\oplus2}\), proving \eqref{eq:galois}.
\end{proof}

The missing factor \(L_K(s,\psi)\) is a Hecke \(L\)-function over a quadratic field, of degree two when viewed over \(\Q\). Its analytic continuation is known; the new issue is its critical-strip zeros. Under the repository's results for \(\zeta\) and \(L(s,\chi_{-23})\), RH for \(L_K(s,\psi)\), RH for \(\zeta_F\), and RH for \(\zeta_H\) are equivalent. Inside the strip all factors in \eqref{eq:cubic}--\eqref{eq:galois} are holomorphic, so zeros cannot be canceled by poles.

This is a better first extension target than another quadratic field, whose zeta function is already a product of two Dirichlet factors. It is also narrower and more diagnostic than attempting all automorphic GRH at once.

\subsection{The local factors are explicit}

Set \(t=p^{-s}\), and write \(D_F(t)\) and \(D_\psi(t)\) for the reciprocal Euler factors of \(\zeta_F\) and \(L_K(\psi)\). The possibilities, including the ramified prime, are:

\begin{center}
\begin{tabular}{llll}
\toprule
Cubic splitting type & \(D_F(t)\) & \(D_\psi(t)\) & Trace at \(p\)\\
\midrule
\(1+1+1\) & \((1-t)^3\) & \((1-t)^2\) & \(2\)\\
\(1+2\) & \((1-t)(1-t^2)\) & \(1-t^2\) & \(0\)\\
\(3\) & \(1-t^3\) & \(1+t+t^2\) & \(-1\)\\
\(p=23\) & \((1-t)^2\) & \(1-t\) & \(1\)\\
\bottomrule
\end{tabular}
\end{center}

Each row satisfies \(D_F(t)=(1-t)D_\psi(t)\). At the ramified prime, the prime ideal of \(K\) has class of order dividing \(2\), hence is principal because the class group has order \(3\). Its \(\psi\)-value is therefore \(1\).

The signed traces show why replacing a Dirichlet character by this new sequence inside a positivity argument is unjustified. There is, however, an exact way to assemble a positive statistic by taking the full class-character product.

\subsection{An exact principal-ideal target}\label{sec:principal-target}

For a nonzero integral ideal \(\mathfrak a\) of \(K\), put
\[
\Lambda_K(\mathfrak a)=
\begin{cases}
\log\Nm\mathfrak p,&\mathfrak a=\mathfrak p^r,\ r\ge1,\\
0,&\text{otherwise}.
\end{cases}
\]
Character orthogonality on the class group gives
\[
1+\psi(\mathfrak a)+\overline{\psi(\mathfrak a)}
=3\,\mathbf1_{[\mathfrak a]=1}.
\]
Taking logarithmic derivatives of \eqref{eq:galois}, for \(\Rea s>1\),
\begin{equation}\label{eq:class-logder}
-\frac{\zeta_H'}{\zeta_H}(s)
=3\sum_{[\mathfrak a]=1}\frac{\Lambda_K(\mathfrak a)}{(\Nm\mathfrak a)^s}.
\end{equation}
All coefficients are nonnegative. The summation is over \emph{principal ideal powers}, not merely powers of principal prime ideals: a nonprincipal prime ideal \(\mathfrak p\) contributes through \(\mathfrak p^3\).

Define the Euler--Kronecker constant
\[
\gamma_H=\lim_{s\to1}\left(\frac{\zeta_H'}{\zeta_H}(s)+\frac1{s-1}\right),
\]
and the normalized smoothed error
\begin{equation}\label{eq:RHprincipal}
\mathcal R_H(x)=
3\!\!\sum_{\substack{\Nm\mathfrak a\le x\\{}[\mathfrak a]=1}}
\Lambda_K(\mathfrak a)
\left(\frac1{\Nm\mathfrak a}-\frac1x\right)
-\log x+(1+\gamma_H).
\end{equation}

\begin{theorem}[A precise sufficient estimate for the new target]\label{thm:principal}
If, for every \(\varepsilon>0\),
\begin{equation}\label{eq:missing}
\mathcal R_H(x)\le x^{-1/2+\varepsilon}
\qquad\text{for all sufficiently large real \(x\),}
\end{equation}
then RH holds for \(L_K(s,\psi)\), \(\zeta_F(s)\), and \(\zeta_H(s)\).
\end{theorem}

\statusbox{\textbf{The unresolved task is \eqref{eq:missing}.} The proof below establishes the analytic reduction to that estimate. Neither the repository's exported Dirichlet theorem nor the finite computations in this archive prove the estimate.}

\begin{proof}
We prove the stronger fixed-boundary implication: for \(0<\beta<1\), an eventual bound \(\mathcal R_H(x)\le x^{\beta-1}\) excludes zeros of \(\zeta_H\) to the right of \(\beta\).

Write \(S_H(x)\) for the sum in \eqref{eq:RHprincipal}. For \(\Rea s>1\), absolute convergence and the elementary integral
\[
\int_n^\infty\left(\frac1n-\frac1x\right)x^{-s}\,dx
=\frac{n^{-s}}{s(s-1)}
\]
give
\begin{align}\label{eq:Phi}
\Phi_H(s)
&:=\int_1^\infty\mathcal R_H(x)x^{-s}\,dx\\
&=
\frac{-\zeta_H'(s)/\zeta_H(s)}{s(s-1)}
-\frac1{(s-1)^2}
+\frac{1+\gamma_H}{s-1}.\nonumber
\end{align}
To check the normalization, put \(u=s-1\). The Laurent expansion is
\[
-\frac{\zeta_H'}{\zeta_H}(1+u)
=\frac1u-\gamma_H+O(u),
\]
and division by \(s(s-1)=u(1+u)\) gives
\[
\frac1{u^2}-\frac{1+\gamma_H}{u}+O(1).
\]
Both singular terms cancel in \eqref{eq:Phi}; \(s=1\) is removable.

At a zero \(\rho\ne0,1\) of multiplicity \(m_\rho\), the remaining transform has a simple pole with residue
\begin{equation}\label{eq:positive-residue}
\frac{m_\rho}{\rho(1-\rho)}.
\end{equation}
Its real residues in \((\beta,1)\) are nonnegative, and there are no real zeros above \(1\), by the Euler product. At other real points to the right of \(\beta\) the residue is zero.

For completeness, the growth and measurability requirements of the general Landau interface are harmless here. Grouping coefficients of \(-\zeta_H'/\zeta_H\) by rational norm gives
\[
0\le\Lambda_H(n)\le[H:\Q]\Lambda(n)=6\Lambda(n).
\]
For \(n=p^k\), this follows by summing the residue degrees of primes above \(p\) whose residue degree divides \(k\); their sum is at most the degree. Thus \(S_H(x)=O(1+\log^2x)\) even using only \(\Lambda(n)\le\log n\). Consequently \(E(v):=\mathcal R_H(e^v)\) is measurable and \(|E(v)|\le Ce^v\) for \(v\ge0\). Dedekind-zeta continuation makes \(\Phi_H\) meromorphic, with at most the simple poles \eqref{eq:positive-residue} in \(\Rea s>0\). The transform identity holds in \(\Rea s>1\), more than the source interface requires.

The assumed bound gives \(E(v)\le e^{(\beta-1)v}\) eventually. Apply the general \code{IsPolar.landau} theorem from \src{Lib/Analytic/Landau.lean}. Every residue with \(\Rea\rho>\beta\) is zero. Equation~\eqref{eq:positive-residue} then excludes every zero there.

Apply this implication for all \(\beta>1/2\). The functional equation gives RH for \(\zeta_H\). The factorization \eqref{eq:galois}, with all factors holomorphic in the strip, then forces the zero-location statement for \(L_K(s,\psi)\), and \eqref{eq:cubic} gives it for \(\zeta_F\).
\end{proof}

The same analytic proof works for any fixed Dedekind zeta function, with its own Euler--Kronecker constant. Its utility here is to identify exactly what a class-group-sensitive arithmetic construction must bound.

This formulation removes the need for \emph{analytic zero-avoiding twists}: the full product has only positive zero multiplicities. Twists may still be necessary inside the arithmetic construction to control denominators and determinant heights. These are two different roles and should not be conflated.

\subsection{The geometry cannot be replaced indiscriminately}

Although \(H\) contains the imaginary quadratic field \(K\) and is totally imaginary, \(H\) is not a CM field. In a Galois CM field, complex conjugation is central; in \(S_3\) it is a noncentral transposition. Replacing the repository's auxiliary CM field by \(H\) and retaining its unitary constructions is therefore invalid.

A more plausible approach retains the CM base \(K\), uses its ideal classes to index suitable theta data, and seeks a determinant whose finite-place divisibility detects the principal-ideal statistic \eqref{eq:RHprincipal}. The formal resemblance is meaningful, but the needed integrality, spanning, height, and main-term identities remain to be proved.
